%%%%%%%%%%%%%%%%%%%%%%%%%%%%%%%%%%%%%%%%%%%%%%%%%%%%%%%%%%%%%%%%%%%
% Physik IV - Blatt 1: Abgabe, Tafelvortrag und Lernloesung
% Standalone one-file LaTeX document
%%%%%%%%%%%%%%%%%%%%%%%%%%%%%%%%%%%%%%%%%%%%%%%%%%%%%%%%%%%%%%%%%%%

\documentclass[11pt,a4paper]{scrartcl}

%%%%%%%%%%%%%%%%%%%%%%%%%%%%%%%%%%%%%%%%%%%%%%%%%%%%%%%%%%%%%%%%%%%
% Embedded file: includes/university_metadata.tex
%%%%%%%%%%%%%%%%%%%%%%%%%%%%%%%%%%%%%%%%%%%%%%%%%%%%%%%%%%%%%%%%%%%
\newcommand{\CourseTitle}{Physik IV: Kerne und Teilchen}
\newcommand{\CourseShortTitle}{Physik IV}
\newcommand{\DocumentKind}{Hausaufgabe / Lernloesung}
\newcommand{\DocumentTitle}{Aufgabenblatt 1}
\newcommand{\DocumentSubtitle}{Einheiten, relativistische kinetische Energie und Streuung an harter Kugel}
\newcommand{\AuthorName}{Tobias Laser}
\newcommand{\InstitutionName}{Universitaet Innsbruck}
\newcommand{\SubmissionDate}{\today}

\newcommand{\makeuniversitytitle}{%
    \begin{titlepage}
        \centering
        \vspace*{2cm}
        {\Large \CourseTitle\par}
        \vspace{1.2cm}
        {\huge\bfseries \DocumentTitle\par}
        \ifx\DocumentSubtitle\empty\else
            \vspace{0.4cm}
            {\Large \DocumentSubtitle\par}
        \fi
        \vspace{1cm}
        {\Large \DocumentKind\par}
        \vfill
        {\large \AuthorName\par}
        \vspace{0.2cm}
        {\large \InstitutionName\par}
        \vspace{0.8cm}
        {\large \SubmissionDate\par}
    \end{titlepage}
    \pagenumbering{arabic}
}

\newcommand{\makechaptertitle}{%
    \begin{center}
        {\Large \CourseTitle\par}
        \vspace{0.5cm}
        {\huge\bfseries \DocumentTitle\par}
        \ifx\DocumentSubtitle\empty\else
            \vspace{0.25cm}
            {\Large \DocumentSubtitle\par}
        \fi
        \vspace{0.5cm}
        {\large \AuthorName \hfill \SubmissionDate\par}
    \end{center}
    \vspace{0.5cm}
}
\newcommand{\makephysiktitle}{\makeuniversitytitle}

%%%%%%%%%%%%%%%%%%%%%%%%%%%%%%%%%%%%%%%%%%%%%%%%%%%%%%%%%%%%%%%%%%%
% Embedded file: includes/university_packages.tex
%%%%%%%%%%%%%%%%%%%%%%%%%%%%%%%%%%%%%%%%%%%%%%%%%%%%%%%%%%%%%%%%%%%
\usepackage[margin=2.5cm]{geometry}
\usepackage[onehalfspacing]{setspace}
\usepackage{scrlayer-scrpage}

\usepackage[T1]{fontenc}
\usepackage[utf8]{inputenc}
\usepackage[ngerman]{babel}
\usepackage{lmodern}
\usepackage{microtype}
\usepackage{csquotes}

\usepackage{amsmath,amsthm,amssymb,mathtools}
\usepackage{bm}
\usepackage{icomma}
\usepackage{cancel}

\usepackage{graphicx}
\usepackage{tikz}
\usetikzlibrary{arrows.meta,calc,decorations.pathmorphing,positioning,patterns,angles,quotes}
\usepackage{pgfplots}
\pgfplotsset{compat=1.18}

\usepackage[locale=DE,space-before-unit=true,per-mode=symbol,separate-uncertainty=true]{siunitx}
\usepackage{booktabs,multirow,array,tabularx,longtable}
\usepackage{caption,subcaption}
\usepackage{placeins}
\usepackage{enumitem}
\usepackage{xparse}

\usepackage[most]{tcolorbox}
\tcbuselibrary{breakable,skins,theorems}

\usepackage[breaklinks=true,colorlinks=true,linkcolor=blue,urlcolor=blue,citecolor=blue]{hyperref}
\usepackage[nameinlink,noabbrev]{cleveref}

\clearpairofpagestyles
\ihead{\CourseShortTitle}
\ohead{\DocumentKind}
\cfoot{\pagemark}
\pagestyle{scrheadings}

\setlist[enumerate]{itemsep=0.4em,topsep=0.4em}
\setlist[itemize]{itemsep=0.25em,topsep=0.35em}
\captionsetup{font=small,labelfont=bf}

%%%%%%%%%%%%%%%%%%%%%%%%%%%%%%%%%%%%%%%%%%%%%%%%%%%%%%%%%%%%%%%%%%%
% Embedded file: includes/university_macros.tex
%%%%%%%%%%%%%%%%%%%%%%%%%%%%%%%%%%%%%%%%%%%%%%%%%%%%%%%%%%%%%%%%%%%
\newcommand{\R}{\mathbb{R}}
\newcommand{\N}{\mathbb{N}}
\newcommand{\Z}{\mathbb{Z}}
\newcommand{\Q}{\mathbb{Q}}
\newcommand{\C}{\mathbb{C}}

\newcommand{\set}[1]{\left\{ #1 \right\}}
\newcommand{\abs}[1]{\left| #1 \right|}
\newcommand{\norm}[1]{\left\lVert #1 \right\rVert}
\newcommand{\avg}[1]{\left\langle #1 \right\rangle}
\newcommand{\paren}[1]{\left( #1 \right)}
\newcommand{\bracket}[1]{\left[ #1 \right]}
\newcommand{\ol}[1]{\overline{#1}}

\newcommand{\dd}{\,\mathrm{d}}
\newcommand{\dv}[2]{\frac{\dd #1}{\dd #2}}
\newcommand{\pdv}[2]{\frac{\partial #1}{\partial #2}}
\newcommand{\pddv}[2]{\frac{\partial^2 #1}{\partial #2^2}}
\newcommand{\evalat}[2]{\left. #1 \right|_{#2}}

\newcommand{\vect}[1]{\bm{#1}}
\newcommand{\unitvect}[1]{\hat{\bm{#1}}}
\newcommand{\grad}{\vect{\nabla}}
\newcommand{\divergence}{\grad\!\cdot}
\newcommand{\curl}{\grad\!\times}

\newcommand{\half}{\frac{1}{2}}
\newcommand{\third}{\frac{1}{3}}
\newcommand{\twothirds}{\frac{2}{3}}
\newcommand{\hbarc}{\hbar c}
\newcommand{\alphaf}{\alpha}
\newcommand{\eV}{\si{eV}}
\newcommand{\MeV}{\si{MeV}}
\newcommand{\GeV}{\si{GeV}}
\newcommand{\TeV}{\si{TeV}}

\newcommand{\nuclide}[3]{^{#1}_{#2}\mathrm{#3}}
\newcommand{\isotope}[2]{^{#1}\mathrm{#2}}
\newcommand{\anti}[1]{\overline{#1}}
\newcommand{\pbar}{\anti{p}}
\newcommand{\nbar}{\anti{n}}
\newcommand{\ee}{e^-}
\newcommand{\ep}{e^+}
\newcommand{\mup}{\mu^+}
\newcommand{\mum}{\mu^-}
\newcommand{\nue}{\nu_e}
\newcommand{\numu}{\nu_\mu}
\newcommand{\nutau}{\nu_\tau}
\newcommand{\pionp}{\pi^+}
\newcommand{\pionm}{\pi^-}
\newcommand{\pionz}{\pi^0}
\newcommand{\kaonp}{K^+}
\newcommand{\kaonm}{K^-}
\newcommand{\kaonz}{K^0}

\newcommand{\diffxs}{\frac{\dd\sigma}{\dd\Omega}}
\newcommand{\totalxs}{\sigma_{\mathrm{total}}}
\newcommand{\lum}{\mathcal{L}}
\newcommand{\smand}{s}
\newcommand{\tmand}{t}
\newcommand{\umand}{u}
\newcommand{\Ecm}{E_{\mathrm{CM}}}

\newcommand{\ie}{d.\,h.}
\newcommand{\eg}{z.\,B.}
\newcommand{\wrt}{bezüglich}

\DeclareSIUnit{\barn}{b}
\DeclareSIUnit{\millibarn}{mb}
\DeclareSIUnit{\microbarn}{\micro b}
\DeclareSIUnit{\nanobarn}{nb}
\DeclareSIUnit{\fermi}{fm}
\DeclareSIUnit{\year}{a}

%%%%%%%%%%%%%%%%%%%%%%%%%%%%%%%%%%%%%%%%%%%%%%%%%%%%%%%%%%%%%%%%%%%
% Embedded file: includes/university_boxes.tex
%%%%%%%%%%%%%%%%%%%%%%%%%%%%%%%%%%%%%%%%%%%%%%%%%%%%%%%%%%%%%%%%%%%
\tcbset{
    unibox/.style={
        enhanced,
        breakable,
        sharp corners,
        boxrule=0.6pt,
        left=1.2mm,
        right=1.2mm,
        top=1mm,
        bottom=1mm,
        before skip=0.8em,
        after skip=0.8em,
        fonttitle=\bfseries
    }
}

\newtcolorbox{calculationbox}[1][]{unibox,title=Rechnung,colback=black!2,colframe=black!55,#1}
\newtcolorbox{sidecalcbox}[1][]{unibox,title=Nebenrechnung,colback=black!1,colframe=black!40,#1}
\newtcolorbox{solutionbox}[1][]{unibox,title=Lösung,colback=blue!2,colframe=blue!45!black,#1}
\newtcolorbox{answerbox}[1][]{unibox,title=Antwort,colback=green!3,colframe=green!45!black,#1}
\newtcolorbox{notebox}[1][]{unibox,title=Notiz,colback=yellow!6,colframe=yellow!45!black,#1}
\newtcolorbox{definitionbox}[1][]{unibox,title=Definition,colback=cyan!3,colframe=cyan!45!black,#1}
\newtcolorbox{warningbox}[1][]{unibox,title=Achtung,colback=red!3,colframe=red!55!black,#1}
\newtcolorbox{summarybox}[1][]{unibox,title=Zusammenfassung,colback=violet!3,colframe=violet!50!black,#1}
\newtcolorbox{formulabox}[1][]{unibox,title=Formel,colback=orange!4,colframe=orange!65!black,#1}
\newtcolorbox{taskbox}[1][]{unibox,title=Aufgabe,colback=white,colframe=black!75,#1}
\newtcolorbox{talkbox}[1][]{unibox,title=Tafelvortrag,colback=teal!3,colframe=teal!45!black,#1}
\newtcolorbox{learningbox}[1][]{unibox,title=Lernidee,colback=purple!3,colframe=purple!45!black,#1}

%%%%%%%%%%%%%%%%%%%%%%%%%%%%%%%%%%%%%%%%%%%%%%%%%%%%%%%%%%%%%%%%%%%
% Embedded file: includes/university_theorems.tex
%%%%%%%%%%%%%%%%%%%%%%%%%%%%%%%%%%%%%%%%%%%%%%%%%%%%%%%%%%%%%%%%%%%
\theoremstyle{definition}
\newtheorem{definition}{Definition}[section]
\newtheorem{example}{Beispiel}[section]
\theoremstyle{plain}
\newtheorem{satz}{Satz}[section]
\newtheorem{lemma}{Lemma}[section]
\theoremstyle{remark}
\newtheorem*{remark}{Bemerkung}

%%%%%%%%%%%%%%%%%%%%%%%%%%%%%%%%%%%%%%%%%%%%%%%%%%%%%%%%%%%%%%%%%%%
\begin{document}

\makeuniversitytitle
\tableofcontents
\newpage

\section*{Vorbemerkung}
\addcontentsline{toc}{section}{Vorbemerkung}
Diese Ausarbeitung ist so strukturiert, dass sie gleichzeitig als saubere Abgabe, als Vorbereitung für einen kurzen Tafelvortrag und als ausführliche Lernunterlage verwendet werden kann. Die Boxen \emph{Tafelvortrag} enthalten jeweils eine knappe, gut abschreibbare Version. Die Boxen \emph{Lernidee} erklären, warum die Rechenschritte physikalisch sinnvoll sind.

\begin{summarybox}[title=Zentrale Themen von Blatt 1]
\begin{itemize}
    \item Umrechnung zwischen Teilchenphysik-Einheiten und SI-Einheiten.
    \item Herleitung der relativistischen kinetischen Energie aus Arbeit und Impuls.
    \item Geometrische Herleitung des Wirkungsquerschnitts für elastische Streuung an einer harten Kugel.
\end{itemize}
\end{summarybox}

\newpage

\section{Aufgabe 1: Einheiten der Kern- und Teilchenphysik}

\begin{taskbox}
\begin{enumerate}[label=(\alph*)]
    \item Berechnen Sie die einem Elektronenvolt entsprechende Energie in SI-Einheiten.
    \item Bestimmen Sie die Ruhemassen folgender Teilchen in SI-Einheiten:
    \begin{enumerate}[label=(\roman*)]
        \item Myon: $m_\mu = \SI{105.6583755}{\MeV}/c^2$,
        \item Neutron: $m_n = \SI{939.56542052}{\MeV}/c^2$,
        \item Tritium: $m_t = \SI{2808.92113668}{\MeV}/c^2$.
    \end{enumerate}
\end{enumerate}
\end{taskbox}

\subsection{Lernidee}

\begin{learningbox}
Das Elektronenvolt ist keine SI-Einheit, sondern eine praktische Energieeinheit. Es ist definiert als die Energie, die eine Ladung $e$ beim Durchlaufen einer Spannung von $\SI{1}{V}$ gewinnt:
\[
    E = qU.
\]
Da $\SI{1}{V}=\SI{1}{J/C}$ gilt, folgt sofort $\SI{1}{eV}=e\cdot\SI{1}{V}$ in Joule. Für Massenangaben der Form $\si{eV}/c^2$ nutzt man anschließend Einsteins Beziehung $E=mc^2$, also $m=E/c^2$.
\end{learningbox}

\subsection{Abgabe-Lösung}

\begin{enumerate}[label=(\alph*)]
    \item \textbf{Ein Elektronenvolt in SI-Einheiten.}

    \begin{calculationbox}
    Mit der Elementarladung
    \[
        e = \SI{1.602176634e-19}{C}
    \]
    und $\SI{1}{V}=\SI{1}{J/C}$ ist
    \[
        \SI{1}{eV}
        = e\cdot \SI{1}{V}
        = \SI{1.602176634e-19}{C}\,\SI{1}{J/C}
        = \SI{1.602176634e-19}{J}.
    \]
    \end{calculationbox}

    \begin{answerbox}
    \[
        \boxed{\SI{1}{eV}=\SI{1.602176634e-19}{J}}.
    \]
    \end{answerbox}

    \item \textbf{Umrechnung von $\si{MeV}/c^2$ in Kilogramm.}

    \begin{formulabox}
    Für eine Masse, die als $M\,\si{MeV}/c^2$ angegeben ist, gilt
    \[
        m_{\mathrm{SI} }
        = M\,\frac{\SI{1}{MeV}}{c^2}
        = M\,\frac{10^6\,\SI{1.602176634e-19}{J}}{(\SI{299792458}{m/s})^2}.
    \]
    \end{formulabox}

    \begin{sidecalcbox}[title=Umrechnungsfaktor]
    \[
        \frac{\SI{1}{MeV}}{c^2}
        = \frac{10^6\cdot \SI{1.602176634e-19}{J}}{(\SI{299792458}{m/s})^2}
        = \SI{1.7826619216e-30}{kg}.
    \]
    \end{sidecalcbox}

    \begin{calculationbox}
    \begin{align*}
        m_\mu
        &= 105.6583755\cdot\SI{1.7826619216e-30}{kg}
         = \SI{1.883531627e-28}{kg},\\
        m_n
        &= 939.56542052\cdot\SI{1.7826619216e-30}{kg}
         = \SI{1.674927498e-27}{kg},\\
        m_t
        &= 2808.92113668\cdot\SI{1.7826619216e-30}{kg}
         = \SI{5.007356751e-27}{kg}.
    \end{align*}
    \end{calculationbox}

    \begin{answerbox}
    \begin{align*}
        \boxed{m_\mu = \SI{1.883531627e-28}{kg}},\\
        \boxed{m_n = \SI{1.674927498e-27}{kg}},\\
        \boxed{m_t = \SI{5.007356751e-27}{kg}}.
    \end{align*}
    \end{answerbox}
\end{enumerate}

\subsection{Tafelvortrag zu Aufgabe 1}
\begin{talkbox}
\begin{align*}
    \SI{1}{eV} &= e\cdot \SI{1}{V}
    = \SI{1.602176634e-19}{C}\,\SI{1}{J/C}
    = \SI{1.602176634e-19}{J},\\
    \frac{\SI{1}{MeV}}{c^2}
    &= \frac{10^6\cdot \SI{1.602176634e-19}{J}}{(\SI{299792458}{m/s})^2}
    = \SI{1.7826619216e-30}{kg}.
\end{align*}
Damit:
\[
    m_\mu=\SI{1.8835e-28}{kg},\qquad
    m_n=\SI{1.6749e-27}{kg},\qquad
    m_t=\SI{5.0074e-27}{kg}.
\]
\end{talkbox}

\FloatBarrier
\newpage

\section{Aufgabe 2: Relativistische Formel der kinetischen Energie}

\begin{taskbox}
Ausgehend von der Arbeit
\[
    W=\int \vect F\cdot\dd\vect r
\]
und in eindimensionaler Betrachtung
\[
    E_{\mathrm{kin}}=\int_0^r F\dd r
\]
soll mit
\[
    p=\gamma mv,\qquad \gamma=\frac{1}{\sqrt{1-v^2/c^2}}
\]
die relativistische kinetische Energie hergeleitet werden.
\end{taskbox}

\subsection{Lernidee}

\begin{learningbox}
Klassisch gilt $F=ma$, aber relativistisch ist die sauberere Definition
\[
    F = \dv{p}{t}.
\]
Der Impuls ist nicht mehr $mv$, sondern $p=\gamma mv$. Deshalb darf man nicht einfach $\int ma\dd r$ rechnen. Stattdessen schreibt man
\[
    \dd r = v\dd t
\]
und erhält dadurch ein Integral über die Geschwindigkeit. Die Energie wächst dann wegen des Lorentzfaktors stärker als klassisch.
\end{learningbox}

\subsection{Abgabe-Lösung}

\begin{enumerate}[label=(\alph*)]
    \item \textbf{Einsetzen von Kraft und relativistischem Impuls.}

    \begin{calculationbox}
    In einer Dimension ist
    \[
        F = \dv{p}{t}.
    \]
    Mit
    \[
        p=\gamma mv=\frac{mv}{\sqrt{1-v^2/c^2}}
    \]
    folgt direkt
    \[
        F
        = \dv{}{t}\paren{\frac{mv}{\sqrt{1-v^2/c^2}}}.
    \]
    Einsetzen in die Arbeitsdefinition liefert
    \[
        E_{\mathrm{kin}}
        =\int_0^r F\dd r
        =\int_0^r
        \dv{}{t}\paren{\frac{mv}{\sqrt{1-v^2/c^2}}}\dd r.
    \]
    \end{calculationbox}

    \begin{answerbox}
    \[
        \boxed{
        E_{\mathrm{kin}}
        =\int_0^r
        \dv{}{t}\paren{\frac{mv}{\sqrt{1-v^2/c^2}}}\dd r
        }.
    \]
    \end{answerbox}

    \item \textbf{Umformung auf ein Integral über $v$ und Lösung.}

    \begin{calculationbox}
    Wegen $v=\dv{r}{t}$ gilt $\dd r=v\dd t$. Damit wird
    \[
        E_{\mathrm{kin}}
        = \int_0^t v\,\dv{p}{t'}\dd t'.
    \]
    Wenn die Geschwindigkeit monoton von $0$ bis $v$ zunimmt, kann man $p$ als Funktion von $v'$ betrachten:
    \[
        E_{\mathrm{kin} }
        = \int_0^v v'\dv{p}{v'}\dd v'.
    \]
    Partielle Integration mit $u=v'$ und $\dd w=\dv{p}{v'}\dd v'=\dd p$ ergibt
    \[
        \int_0^v v'\dd p
        = \evalat{v'p(v')}{0}^{v}-\int_0^v p(v')\dd v'.
    \]
    Mit $p(v')=m v'/\sqrt{1-v'^2/c^2}$ folgt
    \begin{align*}
        E_{\mathrm{kin}}
        &= \frac{mv^2}{\sqrt{1-v^2/c^2}}
           -m\int_0^v \frac{v'}{\sqrt{1-v'^2/c^2}}\dd v'.
    \end{align*}
    Das verbleibende Integral ist
    \begin{align*}
        \int \frac{v'}{\sqrt{1-v'^2/c^2}}\dd v'
        &= -c^2\sqrt{1-v'^2/c^2}.
    \end{align*}
    Deshalb
    \begin{align*}
        E_{\mathrm{kin}}
        &= \frac{mv^2}{\sqrt{1-v^2/c^2}}
           +mc^2\sqrt{1-v^2/c^2}-mc^2.
    \end{align*}
    Mit $\gamma=1/\sqrt{1-v^2/c^2}$ und $\sqrt{1-v^2/c^2}=1/\gamma$ wird
    \begin{align*}
        E_{\mathrm{kin}}
        &= \gamma mv^2 + \frac{mc^2}{\gamma}-mc^2.
    \end{align*}
    Da $v^2=c^2(1-1/\gamma^2)$ gilt,
    \begin{align*}
        \gamma mv^2
        &= mc^2\paren{\gamma-\frac{1}{\gamma}},\\
        E_{\mathrm{kin}}
        &= mc^2\paren{\gamma-\frac{1}{\gamma}}+\frac{mc^2}{\gamma}-mc^2
         = mc^2(\gamma-1).
    \end{align*}
    \end{calculationbox}

    \begin{answerbox}
    \[
        \boxed{E_{\mathrm{kin}} = mc^2(\gamma-1)
        = mc^2\paren{\frac{1}{\sqrt{1-v^2/c^2}}-1}}.
    \]
    \end{answerbox}

    \item \textbf{Reihenentwicklung und klassischer Grenzfall.}

    \begin{calculationbox}
    Für $\beta=v/c$ gilt
    \[
        \gamma = (1-\beta^2)^{-1/2}.
    \]
    Die Binomialreihe liefert für $\abs{\beta}<1$
    \[
        (1-\beta^2)^{-1/2}
        = 1+\frac{1}{2}\beta^2+\frac{3}{8}\beta^4+\mathcal{O}(\beta^6).
    \]
    Damit
    \begin{align*}
        E_{\mathrm{kin}}
        &= mc^2(\gamma-1)\\
        &= mc^2\paren{\frac{1}{2}\beta^2+\frac{3}{8}\beta^4+\mathcal{O}(\beta^6)}\\
        &= \frac{1}{2}mv^2+\frac{3}{8}\frac{mv^4}{c^2}+\mathcal{O}\paren{\frac{v^6}{c^4}}.
    \end{align*}
    \end{calculationbox}

    \begin{answerbox}
    \[
        \boxed{E_{\mathrm{kin}}
        =\frac{1}{2}mv^2+\frac{3}{8}\frac{mv^4}{c^2}+\cdots}
    \]
    Für $v\ll c$ ist der zweite Term sehr klein. Daher bleibt als führender Term das klassische Ergebnis
    \[
        \boxed{E_{\mathrm{kin,klassisch}}=\frac{1}{2}mv^2}.
    \]
    \end{answerbox}
\end{enumerate}

\subsection{Tafelvortrag zu Aufgabe 2}
\begin{talkbox}
\begin{align*}
    E_{\mathrm{kin} }
    &=\int_0^r F\dd r
      =\int_0^r \dv{p}{t}\dd r,
      \qquad
      p=\gamma mv=\frac{mv}{\sqrt{1-v^2/c^2}},\\
    \dd r&=v\dd t
    \quad\Rightarrow\quad
    E_{\mathrm{kin}}=\int_0^v v\dv{p}{v}\dd v.
\end{align*}
Partielle Integration:
\begin{align*}
    E_{\mathrm{kin} }
    &= \evalat{vp}{0}^{v}-\int_0^v p\dd v
     = \frac{mv^2}{\sqrt{1-v^2/c^2}}
       -m\int_0^v\frac{v'}{\sqrt{1-v'^2/c^2}}\dd v'\\
    &= \frac{mv^2}{\sqrt{1-v^2/c^2}}
       +mc^2\sqrt{1-v^2/c^2}-mc^2
     = mc^2(\gamma-1).
\end{align*}
Reihe:
\[
    \gamma=1+\frac12\frac{v^2}{c^2}+\frac38\frac{v^4}{c^4}+\cdots
    \quad\Rightarrow\quad
    E_{\mathrm{kin}}=\frac12mv^2+\frac38\frac{mv^4}{c^2}+\cdots.
\]
\end{talkbox}

\FloatBarrier
\newpage

\section{Aufgabe 3: Streuung an harter Kugel}

\begin{taskbox}
Für die elastische Streuung an einer harten Kugel gilt allgemein
\[
    \diffxs = \frac{b}{\sin\theta}\abs{\dv{b}{\theta}}.
\]
Hierbei ist $b$ der Stoßparameter und $\theta$ der Streuwinkel. Gesucht sind:
\begin{enumerate}[label=(\alph*)]
    \item der Zusammenhang zwischen $b$ und $\theta$ sowie $\dd b$,
    \item der differentielle Wirkungsquerschnitt,
    \item der totale Wirkungsquerschnitt,
    \item der Wirkungsquerschnitt für $40^\circ\leq\theta\leq60^\circ$.
\end{enumerate}
\end{taskbox}

\subsection{Geometrie und Lernidee}

\begin{learningbox}
Der Stoßparameter $b$ beschreibt, wie weit die einlaufende Bahn seitlich am Mittelpunkt vorbeigeht. Für eine harte Kugel findet eine spiegelnde Reflexion an der Kugeloberfläche statt. Kleine $b$ entsprechen fast frontaler Streuung und damit großen Ablenkwinkeln; $b\approx R$ entspricht streifender Streuung und damit $\theta\approx0$.
\end{learningbox}

\begin{figure}[h]
\centering
\begin{tikzpicture}[scale=1.05, >=Latex]
    \coordinate (O) at (0,0);
    \coordinate (P) at (-1.05,1.7);
    \draw[thick] (O) circle (2);
    \draw[->,thick] (-5,1.7) -- (P) node[midway,above] {einlaufend};
    \draw[->,thick] (P) -- ++(-2.0,1.25) node[midway,above] {gestreut};
    \draw[dashed] (-5,0) -- (0,0);
    \draw[dashed] (-5,1.7) -- (-5,0);
    \draw[<->] (-4.8,0) -- (-4.8,1.7) node[midway,left] {$b$};
    \draw[thick] (O) -- (P) node[midway,right] {$R$};
    \draw[->] (-3.9,1.7) arc[start angle=180,end angle=147,radius=1.1];
    \node at (-3.35,2.15) {$\theta$};
    \fill (O) circle (1.2pt) node[below right] {Zentrum};
    \fill (P) circle (1.2pt) node[above] {Kontaktpunkt};
\end{tikzpicture}
\caption{Geometrische Idee der elastischen Streuung an einer harten Kugel.}
\end{figure}

\subsection{Abgabe-Lösung}

\begin{enumerate}[label=(\alph*)]
    \item \textbf{Zusammenhang zwischen $b$ und $\theta$.}

    \begin{calculationbox}
    Bei der spiegelnden Reflexion an der Kugeloberfläche ist der Streuwinkel durch die Geometrie des Radiusvektors am Kontaktpunkt bestimmt. Für den Grenzfall $b=0$ muss $\theta=\pi$ gelten, und für $b=R$ muss $\theta=0$ gelten. Der passende Zusammenhang ist
    \[
        b=R\cos\paren{\frac{\theta}{2}}.
    \]
    Ableiten nach $\theta$ liefert
    \[
        \dv{b}{\theta}
        = -\frac{R}{2}\sin\paren{\frac{\theta}{2}},
        \qquad
        \dd b = -\frac{R}{2}\sin\paren{\frac{\theta}{2}}\dd\theta.
    \]
    Für den Wirkungsquerschnitt verwendet man den Betrag, weil Flächen positiv sind:
    \[
        \abs{\dv{b}{\theta} }
        =\frac{R}{2}\sin\paren{\frac{\theta}{2}}.
    \]
    \end{calculationbox}

    \begin{answerbox}
    \[
        \boxed{b=R\cos\paren{\frac{\theta}{2}}},
        \qquad
        \boxed{\dd b=-\frac{R}{2}\sin\paren{\frac{\theta}{2}}\dd\theta}.
    \]
    \end{answerbox}

    \item \textbf{Differentieller Wirkungsquerschnitt.}

    \begin{calculationbox}
    Einsetzen in
    \[
        \diffxs = \frac{b}{\sin\theta}\abs{\dv{b}{\theta}}
    \]
    ergibt
    \begin{align*}
        \diffxs
        &= \frac{R\cos(\theta/2)}{\sin\theta}
           \cdot \frac{R}{2}\sin(\theta/2)\\
        &= \frac{R^2}{2}\frac{\sin(\theta/2)\cos(\theta/2)}{\sin\theta}.
    \end{align*}
    Mit
    \[
        \sin\theta = 2\sin\paren{\frac{\theta}{2}}\cos\paren{\frac{\theta}{2}}
    \]
    folgt
    \[
        \diffxs = \frac{R^2}{4}.
    \]
    \end{calculationbox}

    \begin{answerbox}
    \[
        \boxed{\diffxs=\frac{R^2}{4}}.
    \]
    Der differentielle Wirkungsquerschnitt ist unabhängig von $\theta$. Die Streuung ist im Raumwinkel isotrop: pro gleichem Raumwinkelelement $\dd\Omega$ werden gleich viele Teilchen erwartet. Pro gleichem Winkelintervall $\dd\theta$ ist die Häufigkeit wegen $\dd\Omega=2\pi\sin\theta\dd\theta$ aber nicht konstant, sondern proportional zu $\sin\theta$.
    \end{answerbox}

    \item \textbf{Totaler Wirkungsquerschnitt.}

    \begin{calculationbox}
    Der totale Wirkungsquerschnitt ist das Integral über den ganzen Raumwinkel:
    \[
        \sigma_{\mathrm{tot}}
        = \int \diffxs\dd\Omega.
    \]
    Da $\diffxs=R^2/4$ konstant ist und der volle Raumwinkel $4\pi$ beträgt,
    \[
        \sigma_{\mathrm{tot}}
        =\frac{R^2}{4}\cdot4\pi
        =\pi R^2.
    \]
    Alternativ explizit:
    \begin{align*}
        \sigma_{\mathrm{tot}}
        &= \int_0^{2\pi}\int_0^\pi \frac{R^2}{4}\sin\theta\dd\theta\dd\phi\\
        &= \frac{R^2}{4}\cdot2\pi\cdot\bracket{-\cos\theta}_0^\pi
         = \pi R^2.
    \end{align*}
    \end{calculationbox}

    \begin{answerbox}
    \[
        \boxed{\sigma_{\mathrm{tot}}=\pi R^2}.
    \]
    Das ist genau die geometrische Schattenfläche der Kugel. Physikalisch bedeutet das: Alle Teilchen mit $b\le R$ treffen die harte Kugel und werden gestreut; alle Teilchen mit $b>R$ gehen vorbei.
    \end{answerbox}

    \item \textbf{Winkelbereich $40^\circ\leq\theta\leq60^\circ$.}

    \begin{calculationbox}
    Für einen Winkelbereich $\theta_1\le\theta\le\theta_2$ gilt
    \begin{align*}
        \sigma(\theta_1\le\theta\le\theta_2)
        &= \int_0^{2\pi}\int_{\theta_1}^{\theta_2}\frac{R^2}{4}\sin\theta\dd\theta\dd\phi\\
        &= \frac{R^2}{4}\cdot2\pi\cdot\bracket{-\cos\theta}_{\theta_1}^{\theta_2}\\
        &= \frac{\pi R^2}{2}\paren{\cos\theta_1-\cos\theta_2}.
    \end{align*}
    Für $\theta_1=40^\circ$ und $\theta_2=60^\circ$ folgt
    \begin{align*}
        \sigma(40^\circ\le\theta\le60^\circ)
        &= \frac{\pi R^2}{2}\paren{\cos 40^\circ-\cos 60^\circ}\\
        &\approx \frac{\pi R^2}{2}(0{,}7660-0{,}5000)\\
        &\approx 0{,}4178\,R^2.
    \end{align*}
    Bezogen auf den totalen Wirkungsquerschnitt ist der Anteil
    \[
        \frac{\sigma(40^\circ\le\theta\le60^\circ)}{\sigma_{\mathrm{tot}}}
        =\frac{\cos40^\circ-\cos60^\circ}{2}
        \approx 0{,}1330.
    \]
    \end{calculationbox}

    \begin{answerbox}
    \[
        \boxed{\sigma(40^\circ\le\theta\le60^\circ)
        =\frac{\pi R^2}{2}\paren{\cos40^\circ-\cos60^\circ}
        \approx 0{,}4178\,R^2}.
    \]
    Das sind ungefähr $13{,}3\%$ des totalen Wirkungsquerschnitts $\pi R^2$.
    \end{answerbox}
\end{enumerate}

\subsection{Tafelvortrag zu Aufgabe 3}
\begin{talkbox}
Geometrie der harten Kugel:
\[
    b=R\cos\paren{\frac{\theta}{2}},
    \qquad
    \abs{\dv{b}{\theta}}=\frac{R}{2}\sin\paren{\frac{\theta}{2}}.
\]
Damit
\begin{align*}
    \diffxs
    &=\frac{b}{\sin\theta}\abs{\dv{b}{\theta}}\\
    &=\frac{R\cos(\theta/2)}{\sin\theta}\frac{R}{2}\sin(\theta/2)
     =\frac{R^2}{4},
\end{align*}
weil $\sin\theta=2\sin(\theta/2)\cos(\theta/2)$.
Total:
\[
    \sigma_{\mathrm{tot}}
    =\int\frac{R^2}{4}\dd\Omega
    =\frac{R^2}{4}\cdot4\pi
    =\pi R^2.
\]
Winkelbereich:
\[
    \sigma_{40^\circ\ldots60^\circ}
    =\frac{R^2}{4}\cdot2\pi\int_{40^\circ}^{60^\circ}\sin\theta\dd\theta
    =\frac{\pi R^2}{2}\paren{\cos40^\circ-\cos60^\circ}.
\]
\end{talkbox}

\FloatBarrier
\newpage

\section{Gesamtzusammenfassung für die Prüfungsvorbereitung}

\begin{summarybox}[title=Was man sicher können sollte]
\begin{enumerate}
    \item \textbf{Einheiten:} $\SI{1}{eV}=\SI{1.602176634e-19}{J}$ und $\SI{1}{MeV}/c^2=\SI{1.7826619216e-30}{kg}$.
    \item \textbf{Relativistische Energie:} Nicht $F=ma$ verwenden, sondern $F=\dd p/\dd t$ mit $p=\gamma mv$.
    \item \textbf{Endformel:} $E_{\mathrm{kin}}=mc^2(\gamma-1)$.
    \item \textbf{Klassischer Grenzfall:} Für $v\ll c$ wird $E_{\mathrm{kin}}\approx \half mv^2$.
    \item \textbf{Harte Kugel:} Aus $b=R\cos(\theta/2)$ folgt $\diffxs=R^2/4$ und $\sigma_{\mathrm{tot}}=\pi R^2$.
\end{enumerate}
\end{summarybox}

\begin{warningbox}[title=Typische Fehler]
\begin{itemize}
    \item Bei $\si{MeV}/c^2$ nicht vergessen: zuerst $\si{MeV}$ in Joule umrechnen, dann durch $c^2$ teilen.
    \item Bei Aufgabe 2 nicht $F=ma$ als Ausgangspunkt benutzen; relativistisch ist $F=\dd p/\dd t$ grundlegend.
    \item Beim Wirkungsquerschnitt den Betrag $\abs{\dd b/\dd\theta}$ verwenden, weil Querschnitte positiv sind.
    \item Beim Winkelbereich in Aufgabe 3 über den Raumwinkel integrieren: $\dd\Omega=\sin\theta\dd\theta\dd\phi$.
\end{itemize}
\end{warningbox}

\end{document}
